\documentclass{ltjsarticle}

\usepackage{luatexja}
\usepackage{amsmath, amssymb}
\usepackage{graphicx}
\usepackage{hyperref}

\title{日本語LaTeXのひな形}
\author{GK5674}
\date{\today}

\begin{document}

\maketitle

\section{はじめに}

これはGitHub Actionsで自動ビルドする日本語LaTeX文書です。

日本語が正しく表示されるか確認します。

\section{数式の例}

二次方程式の解の公式は次の通りです。

\[
x = \frac{-b \pm \sqrt{b^2 - 4ac}}{2a}
\]

\section{箇条書きの例}

\begin{itemize}
  \item 日本語が使える
  \item 数式が使える
  \item GitHub ActionsでPDF化できる
\end{itemize}

\section{まとめ}

このひな形をベースに、レポート・論文・資料を作れます。

\end{document}
